\subsection{Matrices}

The book-proofs in this chapter were a lot harder than previous chapters, but I could still follow them fine.

The problems were \emph{really} easy compared to last subchapter, it makes me wonder what sort of demonic influence that subchapter had\dots

Highlight to \#2, its a great problem!

\begin{enumerate}
    \item[1] Suppose $T\in\mathcal{L}(V,W)$. Show that, with respect to any choices of bases of $V$ and $W$, the matrix of $T$ has at least $\dim\big(\operatorname{range} T\big)$ nonzero entries.
    
    \begin{proof}
        The dimension of the range of $T$ is equivalent to the rank of $\mathcal M(T)$, since matrix-vector multiplication is a linear combination of each column of $\mathcal M(T)$. Therefore, there is some set of $\dim \lrange T$ linearly independent vectors. Each of these vectors must not equal zero, so there are at least $\dim \lrange T$ non-zero entries.
    \end{proof}

    \item[2] Suppose $T \in \mathcal{L}(V, W)$, where $V$ and $W$ are finite-dimensional and nonzero. Prove that $\dim \operatorname{range} T = 1$ if and only if there exist a basis of $V$ and a basis of $W$ such that with respect to these bases, all entries of $\mathcal{M}(T)$ equal 1.
    
    \begin{proof}
        If $\dim \lrange T = 1$, then there is some $w_b \in W$ where every column of $\mathcal M(T)$ is a multiple of $w_b$. 

        We would like to show there is some basis of $W$ such that $w_b = w_1 + \dots + w_m$. This itself is equivalent to showing that there is some linear combination that has all non-zero entries (proof will be presented when returning the other direction). 

        Let us state the lemma we want to prove.

        \begin{lemma*}
            Given vector $v \in V$, $v \neq 0$, there exists a basis $v_1, \dots, v_n$ such that 
            
            $$v = a_1v_1 + \dots + a_nv_n$$
            and
            $$a_i \neq 0, i \in \{1, \dots, n\}$$


        \end{lemma*}

        \begin{proof}
            Choose $v_1 = v$. We then extend this to a basis, $v_1, \dots, v_n$. Because for a given linear combination, each element has only one representation, we know that 

            $$v = v_1 + 0v_2 + \dots + 0v_n$$

            We now define a new basis, $v_1 + \dots + v_n, v_2, \dots, v_n$. It remains a basis (unproved but obvious). Using this basis, 

            $$v = 1(v_1 + \dots + v_n) - 1v_2 - \dots - 1v_n$$

            None of the coefficients are zero, so we have constructed a solution.
        \end{proof}

        Now we can proceed with the main proof.

        If the entries of the matrix are all 1, it is obvious that $\dim \lrange T = 1$. 

        Then, we just need to prove the other way. We are given $T \in \mathcal L(V, W)$, $\dim \lrange T = 1$. 

        Choose some $v_1 \in V$ where $Tv_1 \neq 0$, and let $Tv_1 = w$. By our lemma, we can construct a basis $t_1, \dots, t_n$ where $w = a_1t_1 + \dots + a_nt_n$, $a_i \neq 0$. We define our new basis, then as 

        $$w_i = a_it_i$$

        Now, 

        $$w = w_1 + \dots + w_n$$

        As intended. 

        As for the rest of the basis of $V$, we know that 

        \begin{align*}
            \dim \lrange T + \dim \lnull T &= \dim V \\
            1 + \dim\lnull T &= n \\
            \dim\lnull T &= n - 1
        \end{align*}

        So I can define $v_2, \dots, v_n$ as a basis for $\lnull T$. We know it is LI from $v_1$ obviously. So, we can use 

        $$v_1, v_1 + v_2, \dots, v_n + v_1$$

        as our basis. 

        For any $j \in \{2, \dots, n\}$, $T(v_1 + v_j) = T(v_1) + T(v_j) = w_1 + \dots + w_n$, so our matrix is all ones as desired!
    \end{proof}

    \item[4] Suppose that $D\in\mathcal L\big(\mathcal P_3(\mathbb R),\mathcal P_2(\mathbb R)\big)$ is the differentiation map defined by $Dp=p'$. Find a basis of $\mathcal P_3(\mathbb R)$ and a basis of $\mathcal P_2(\mathbb R)$ such that the matrix of $D$ with respect to these bases is

    \[
    [D] = \begin{pmatrix}
    1 & 0 & 0 & 0 \\
    0 & 1 & 0 & 0 \\
    0 & 0 & 1 & 0
    \end{pmatrix}.
    \]

    \begin{proof}
        \[
        \begin{blockarray}{ccccc}
            & x^3 & x^2 & x & 1 \\
            \begin{block}{c(cccc)}
                3x^2 & 1 & 0 & 0 & 0 \\
                2x & 0 & 1 & 0 & 0 \\
                1 & 0 & 0 & 1 & 0 \\
            \end{block}
        \end{blockarray}
        \]

        I learned blockarray for this!
    \end{proof}

    \item[12] Prove that matrix multiplication is associative. In other words, suppose
    $A$, $B$, and $C$ are matrices of sizes such that $(AB)C$ is defined. Explain why
    $A(BC)$ is defined as well, and prove that

    \[ (AB)C = A(BC). \]

    \begin{proof}
        We can treat $A, B$, and $C$ as linear maps, in which they are just functions.

        Thus, 

        \[(AB)Cv = A(B(C(v))) = A(BC)v\]
    \end{proof}

    \item[13] Suppose $A$ is an $n$-by-$n$ matrix and $1 \le j, k \le n$. Show that the entry in
    row $j$, column $k$, of $A^3$ (which is defined to mean $AAA$) is
    \[
        \sum_{p=1}^n \sum_{r=1}^n A_{j,p} A_{p,r} A_{r,k}.
    \]

    \begin{proof}
        \begin{align*}
            (AA)_{j, k} &= \sum_{r = 1}^n A_{j, r}A_{r, k} \\
            (A(AA))_{j, k} &= \sum_{p=1}^n A_{j, p} (AA)_{p, k} \\
            (A(AA))_{j, k} &= \sum_{p=1}^n A_{j, p} \sum_{r = 1}^n A_{p, r}A_{r, k} \\
            (AAA)_{j, k} &= \sum_{p=1}^n \sum_{r = 1}^n A_{j, p}  A_{p, r}A_{r, k} 
        \end{align*}
        
    \end{proof}

    \item[16] Suppose $A$ is an $m$-by-$n$ matrix with $A \ne 0$. Prove that the rank of $A$ is $1$ if and only if there exist $(c_1, \dots, c_m) \in \mathbb{F}^m$ and $(d_1, \dots, d_n) \in \mathbb{F}^n$ such that
    $A_{j,k} = c_j d_k$
    for every $j = 1, \dots, m$ and every $k = 1, \dots, n$.

    \begin{proof}
        When the rank of $A$ is 1, all of the columns are a multiple of a single non-zero column. We can set $c_1, \dots, c_m$ to be that column, then $d_1, \dots, d_n$ to be the actual multiples. It is clear this is a biconditional.
    \end{proof}

    \item[17] Suppose $T\in\mathcal L(V)$, and $u_1,\dots,u_n$ and $v_1,\dots,v_n$ are bases of $V$. Prove that the following are equivalent:
    \begin{enumerate}
        \item $T$ is injective.
        \item The columns of $\mathcal M(T)$ are linearly independent in $\mathbb F^{n\times 1}$.
        \item The columns of $\mathcal M(T)$ span $\mathbb F^{n\times 1}$.
        \item The rows of $\mathcal M(T)$ span $\mathbb F^{1\times n}$.
        \item The rows of $\mathcal M(T)$ are linearly independent in $\mathbb F^{1\times n}$.
    \end{enumerate}
    Here $\mathcal M(T)$ denotes $\mathcal M\big(T,(u_1,\dots,u_n),(v_1,\dots,v_n)\big)$.

    \begin{proof}
        (a) and (b) are equivalent because any $Tv$ is a linear combination of the columns of $\mathcal M(T)$ using the $v$ for the coefficients. Assuming the columns of $\mathcal M(T)$ are linearly independent, that makes every resultant vector unique, which satisfies the definition of injectivity. If $T$ was injective, then we can show by contradiction that the columns of $\mathcal M(T)$ must be linearly independent.

        (b) and (c) are equivalent because $\dim \F^{n \times 1} = n$, so any length n list of linearly independent vectors will also span the whole space, and vice versa.

        The same reasoning applies to (d) and (e).

        Because row rank and column rank are identical, (c) and (d) are equivalent since the dimensions of both must be $n$ (and therefore they both span the space).

        Combining all of these, we have proved the equivalences.
    \end{proof}

\end{enumerate}